\subsection{Simulation Implementation}
\label{sec:simulation}

The Python simulation implements a two-layer control architecture that accounts for actuator dynamics. Each robot is modeled as an omnidirectional point mass with a velocity controller exhibiting first-order response with time constant $\tau$.

\subsubsection{Actuator Dynamics}

The continuous-time actuator dynamics follow:
\begin{equation}
\dot{\mathbf{v}}(t) = -\frac{1}{\tau}\mathbf{v}(t) + \frac{1}{\tau}\mathbf{v}_{\text{cmd}}(t)
\label{eq:continuous}
\end{equation}

This produces the characteristic exponential response:
\begin{equation}
\mathbf{v}(t) = \mathbf{v}_{\text{cmd}}(1 - e^{-t/\tau})
\end{equation}

For digital implementation at control period $\Delta t$, discretization of Eq.~\eqref{eq:continuous} yields:
\begin{equation}
\mathbf{v}[k+1] = \alpha \mathbf{v}[k] + (1-\alpha)\mathbf{v}_{\text{cmd}}[k]
\label{eq:momentum}
\end{equation}
where the momentum coefficient is $\alpha = e^{-\Delta t/\tau}$.

\subsubsection{Formation Control Integration}

The desired velocity commands are generated by proportional formation control:
\begin{equation}
\mathbf{v}_{\text{des}}[k] = k_v(\mathbf{x}_{\text{des}}[k] - \mathbf{x}[k])
\end{equation}
where $k_v$ is the position error gain and $\mathbf{x}_{\text{des}}[k]$ is the desired robot position computed from the formation kinematics.

Substituting into Eq.~\eqref{eq:momentum} produces the final control law:
\begin{equation}
\mathbf{v}[k+1] = \alpha \mathbf{v}[k] + (1-\alpha)k_v(\mathbf{x}_{\text{des}}[k] - \mathbf{x}[k])
\label{eq:final_control}
\end{equation}

\subsubsection{Control Architecture}

The hierarchical control architecture operates as follows:
\begin{enumerate}
    \item \textbf{High-level controller:} The control primitive computes desired cluster centroid velocity $\mathbf{v}_c^{\text{des}}$ based on field readings
    \item \textbf{Inverse kinematics:} The formation controller uses the inverse Jacobian to map centroid velocity to individual robot desired positions $\mathbf{x}_{\text{des}}[k]$
    \item \textbf{Low-level controller:} Each robot executes Eq.~\eqref{eq:final_control} to track its desired position while accounting for actuator dynamics
\end{enumerate}

This architecture mirrors the physical testbed implementation described in Section~\ref{sec:testbed}.

\subsubsection{Simulation Parameters}

The simulation uses a 10~Hz control loop ($\Delta t = 0.1$~s) with actuator time constant $\tau = 0.3$~s, yielding $\alpha \approx 0.717$. This momentum coefficient characterizes actuator response: $\alpha \to 1$ represents sluggish actuators with high inertia, while $\alpha \to 0$ represents fast, responsive actuators. The position gain is set to $k_v = 1.0$.
